\begin{center}
\SourceLeaf{055}\pdfbookmark[2]{37.}{article-037}37.
\end{center}
Resumamus in motu elliptico aequationem inter anomaliam excentricam et tempus
\[
E-e\sin E=\frac{kt\BareRoot{(1+\mu)}}{a^{\frac{3}{2}}},
\]
ubi $E$ in partibus radii expressam supponimus. Factorem $\BareRoot{(1+\mu)}$ abhinc omittemus; si unquam casus occurreret, ubi eius rationem habere in potestate operaeque pretium esset, signum $t$ non tempus ipsum post perihelium, sed hoc tempus per $\BareRoot{(1+\mu)}$ multiplicatum exprimere deberet. Designamus porro per $q$ distantiam in perihelio, et pro $E$ et $\sin E$ introducimus quantitates $E-\sin E$ et $E-\tfrac{1}{10}(E-\sin E)=\tfrac{9}{10}E+\tfrac{1}{10}\sin E$: rationem cur has potissimum eligamus lector attentus ex sequentibus sponte deprehendet. Hoc modo aequatio nostra formam sequentem induit:
\[
(1-e)(\tfrac{9}{10}E+\tfrac{1}{10}\sin E)+(\tfrac{1}{10}+\tfrac{9}{10}e)(E-\sin E)=kt\left(\frac{1-e}{q}\right)^{\frac{3}{2}}.
\]
Quatenus $E$ ut quantitas parva ordinis primi spectatur, erit $\tfrac{9}{10}E+\tfrac{1}{10}\sin E=E-\tfrac{1}{60}E^3+\tfrac{1}{1200}E^5-\text{etc.}$ quantitas ordinis primi, contra $E-\sin E=\tfrac{1}{6} E^3-\tfrac{1}{120}E^5+\tfrac{1}{5040}E^7-\text{etc.}$ quantitas ordinis tertii. Statuendo itaque
\[
\frac{6(E-\sin E)}{\frac{9}{10}E+\frac{1}{10}\sin E}=4A,\qquad
\frac{\frac{9}{10}E+\frac{1}{10}\sin E}{2\BareRoot{A}}=B,
\]
erit $4A=E^2-\tfrac{1}{30}E^4-\tfrac{1}{5040}E^6-\text{etc.}$ quantitas ordinis secundi, atque $B=1+\tfrac{3}{2800}E^4-\text{etc.}$ ab unitate quantitate quarti ordinis diversa. Aequatio nostra autem hinc fit
\[
\text{1)}\qquad B\left\{2(1-e)A^{\frac{1}{2}}+\tfrac{2}{15}(1+9e)A^{\frac{3}{2}}\right\}=kt\left(\frac{1-e}{q}\right)^{\frac{3}{2}}.
\]

Per tabulas vulgares trigonometricas $\tfrac{9}{10}E+\tfrac{1}{10}\sin E$ quidem praecisione sufficiente calculari potest, non tamen $E-\sin E$, quoties $E$ est angulus parvus: hacce igitur via quantitates $A$ et $B$ satis exacte determinare non liceret. Huic autem difficultati remedium afferret tabula peculiaris, ex qua cum argumento $E$ aut ipsum $B$ aut logarithmum ipsius $B$ excerpere possemus: subsidia ad constructionem talis tabulae necessaria cuique in analysi vel mediocriter versato facile se offerent. Adiumento aequationis
\[
\frac{9E+\sin E}{20B}=\BareRoot{A}
\]
\SourceLeaf{056}etiam $\BareRoot{A}$, atque hinc $t$ per formulam 1), omni quae desiderari potest praecisione determinare liceret.

Ecce specimen talis tabulae, quod saltem léntam augmentationem ipsius $\log B$ manifestabit: superfluum esset, hanc tabulam maiori extensione elaborare, infra enim tabulas formae multo commodioris descripturi sumus:
\begin{center}
\begin{tabular}{|r|r||r|r||r|r|}
\hline
$E$ & $\log B$ & $E$ & $\log B$ & $E$ & $\log B$\\
\hline
$0^0$ & $0{,}000\,0000$ & $25^0$ & $0{,}000\,0168$ & $50^0$ & $0{,}000\,2675$\\
$5$ & $00$ & $30$ & $0349$ & $55$ & $3910$\\
$10$ & $04$ & $35$ & $0645$ & $60$ & $5526$\\
$15$ & $22$ & $40$ & $1099$ & & \\
$20$ & $69$ & $45$ & $1758$ & & \\
\hline
\end{tabular}
\end{center}
